\chapter{Implementation}
\label{ch:implementation}
\section{Code organization}
The Python package under \CodePath{src/transcreation} separates stages for planning, grounding, segmentation, inpainting, appearance transfer, instruction editing and reference replacement. Command-line entry points in \CodePath{scripts} call those stages. The Phase 9 pilot adds a manifest-driven runner with \CodePath{validate}, \CodePath{run}, \CodePath{verify} and \CodePath{package} operations.

\section{Execution environments}
The project targets Windows with Python 3.11. The CPU-only pilot uses Pillow and NumPy; live model stages require pinned model weights, PyTorch and an RTX 3050 Laptop GPU with 4 GB VRAM. Components load one at a time. The recorded Phase 8 SDXL+Canny fallback used FP16 and sequential CPU offload; it is not a FLUX result.

\section{Automated and guided boundaries}
The real automated example reached \texttt{REVIEW} with zero released candidates. Guided experiments carry \texttt{RESEARCH\_ONLY} scope, no approval flag and their own frozen request. This distinction prevents a reviewer from mistaking a manually selected object substitution for a planner-authorized cultural decision.

\section{Verification implementation}
Verification reopens saved PNGs, checks SHA-256 hashes, validates manifest provenance, measures permitted and protected pixel differences, searches for leftover source content in reconstruction regions and performs deterministic in-memory replay. Output directories refuse overwrite. Offline unit tests use synthetic images and mocks; real-model evidence is reported separately.
